\chapter{2014年辽宁卷（理）}

\section{单选题}

\begin{problem}
已知全集 \(U = \R\)，\(A = \{x \mid x \leqslant 0\}\)，\(B = \{x \mid x \geqslant 1\}\)，则集合 \(\complement_U (A \cup B) =\)（\quad）

\choices
    {\( \{x \mid x \geqslant 0\} \)}
    {\(\{x\mid x\leqslant 1\}\)}
    {\(\{x\mid 0\leqslant x\leqslant 1\}\)}
    {\(\{x\mid 0 < x < 1\}\)}
\end{problem}

\begin{answer}
D
\end{answer}

\begin{solution}
因为
\[
A\cup B=\{x\mid x\leqslant0\text{ 或 }x\geqslant1\}
\]
所以
\[
\complement_U(A\cup B)=\{x\mid 0<x<1\}
\]
故选 D.
\end{solution}

\begin{problem}
设复数 \(z\) 满足 \((z - 2\i)(2 - \i) = 5\)，则 \(z =\)（\quad）

\choices
    {\(2 + 3\i\)}
    {\(2 - 3\i\)}
    {\(3 + 2\i\)}
    {\(3 - 2\i\)}
\end{problem}

\begin{answer}
A
\end{answer}

\begin{solution}
由
\[
z-2\i=\frac{5}{2-\i}=\frac{5(2+\i)}{(2-\i)(2+\i)}=2+\i
\]
得 \(z=2+3\i\). 故选 A.
\end{solution}

\begin{problem}
已知 \(a = 2^{-\frac{1}{3}}\)，\(b = \log_{2} \frac{1}{3}\)，\(c = \log_{\frac{1}{2}} \frac{1}{3}\)，则（\quad）

\choices
    {\(a > b > c\)}
    {\(a > c > b\)}
    {\(c > a > b\)}
    {\(c > b > a\)}
\end{problem}

\begin{answer}
C
\end{answer}

\begin{solution}
因为
\(b=\log_2\frac13=-\log_2 3<0\)，\(c=\log_{\frac12}\frac13=\log_2 3>1\)
而 \(0<a=2^{-1/3}<1\)，所以 \(c>a>b\). 故选 C.
\end{solution}

\begin{problem}
已知 \(m\)，\(n\) 表示两条不同直线，\(\alpha\) 表示平面，下列说法正确的是（\quad）

\choices
    {若 \(m \parallel \alpha, n \parallel \alpha\)，则 \(m \parallel n\)}
    {若 \(m \perp \alpha\)，\(n \subset \alpha\)，则 \(m \perp n\)}
    {若 \(m \perp \alpha\)，\(m \perp n\)，则 \(n \parallel \alpha\)}
    {若 \(m \parallel \alpha\)，\(m \perp n\)，则 \(n \perp \alpha\)}
\end{problem}

\begin{answer}
B
\end{answer}

\begin{solution}
若 \(m\perp\alpha\)，则直线 \(m\) 的方向垂直于平面 \(\alpha\) 内任意直线的方向. 因此当 \(n\subset\alpha\) 时，有 \(m\perp n\)，选项 B 正确.

其余选项均不成立：两条都平行于同一平面的直线不一定互相平行；当 \(m\perp\alpha\) 且 \(m\perp n\) 时，\(n\) 也可能是平面 \(\alpha\) 内经过垂足的直线；当 \(m\parallel\alpha\) 且 \(m\perp n\) 时，\(n\) 不一定垂直于 \(\alpha\). 故选 B.
\end{solution}

\begin{problem}
设 \(a\)，\(b\)，\(c\) 是非零向量，已知命题 \(p\)：若 \(a \cdot b = 0\)，\(b \cdot c = 0\)，则 \(a \cdot c = 0\)；命题 \(q\)：若 \(a \parallel b\)，\(b \parallel c\)，则 \(a \parallel c\)，则下列命题中的真命题是（\quad）

\choices
    {\(p \vee q\)}
    {\(p \wedge q\)}
    {\((\neg p) \land (\neg q)\)}
    {\(p \vee (\neg q)\)}
\end{problem}

\begin{answer}
A
\end{answer}

\begin{solution}
命题 \(p\) 不一定成立. 例如取
\(a=c=(1,0,0)\)，\(b=(0,0,1)\)
则 \(a\cdot b=b\cdot c=0\)，但 \(a\cdot c=1\neq0\)，所以 \(p\) 为假.

非零向量的平行关系具有传递性，因此由 \(a\parallel b\) 和 \(b\parallel c\) 可得 \(a\parallel c\)，命题 \(q\) 为真. 所以 \(p\vee q\) 为真，故选 A.
\end{solution}

\begin{problem}
6 把椅子摆成一排，3 人随机就座，任何两人不相邻的坐法种数为（\quad）

\choices
    {144}
    {120}
    {72}
    {24}
\end{problem}

\begin{answer}
D
\end{answer}

\begin{solution}
先从 6 把椅子中选出 3 把互不相邻的椅子. 设选中的椅子编号为 \(i_1<i_2<i_3\)，则 \(i_2-i_1\geqslant2\)，\(i_3-i_2\geqslant2\). 令
\(j_1=i_1\)，\(j_2=i_2-1\)，\(j_3=i_3-2\)
则 \(1\leqslant j_1<j_2<j_3\leqslant4\)，所以选椅子的方式有 \(\mathrm C_4^3=4\) 种.

3 人就座的排列有 \(3!=6\) 种，故坐法种数为
\[
\mathrm C_4^3\cdot3!=4\cdot6=24
\]
故选 D.
\end{solution}

\begin{problem}
某几何体三视图如图所示，则该几何体的体积为（\quad）

\bitmapfigure[max width=0.56\linewidth]{img/2014/liaoning_science/q7_fig1.png}

\choices
    {\(8 - 2\pi\)}
    {\(8 - \pi\)}
    {\(8 - \frac{\pi}{2}\)}
    {\(8 - \frac{\pi}{4}\)}
\end{problem}

\begin{answer}
B
\end{answer}

\begin{solution}
由主视图和左视图可知，几何体的高为 \(2\)，俯视图中的外框为边长 \(2\) 的正方形. 俯视图中被去掉的两块均为半径 \(1\) 的四分之一圆，故底面积为
\[
S=2\times2-2\times\frac14\pi\times1^2=4-\frac\pi2
\]
因此几何体的体积为
\[
V=2S=2\left(4-\frac\pi2\right)=8-\pi
\]
故选 B.
\end{solution}

\begin{problem}
设等差数列 \(\{a_{n}\}\) 的公差为 \(d\)，若数列 \(\{2^{a_{1}a_{n}}\}\) 为递减数列，则（\quad）

\choices
    {\(d < 0\)}
    {\(d > 0\)}
    {\(a_{1}d > 0\)}
    {\(a_{1}d < 0\)}
\end{problem}

\begin{answer}
D
\end{answer}

\begin{solution}
由等差数列的通项公式，\(a_{n+1}-a_n=d\). 由于底数 \(2>1\)，数列 \(\{2^{a_1a_n}\}\) 递减等价于指数数列 \(\{a_1a_n\}\) 递减，即
\[
a_1a_{n+1}<a_1a_n
\]
两边相减得 \(a_1(a_{n+1}-a_n)=a_1d<0\). 故选 D.
\end{solution}

\begin{problem}
将函数 \( y = 3\sin \left(2x + \frac{\pi}{3}\right) \) 的图象向右平移 \( \frac{\pi}{2} \) 个单位长度. 所得图象对应的函数（\quad）

\choices
    {在区间 \( \left[\frac{\pi}{12}, \frac{7\pi}{12}\right] \) 上单调递减}
    {在区间 \(\left[\frac{\pi}{12}, \frac{7\pi}{12}\right]\) 上单调递增}
    {在区间 \( \left[-\frac{\pi}{6}, \frac{\pi}{3}\right] \) 上单调递减}
    {在区间 \( \left[-\frac{\pi}{6},\frac{\pi}{3}\right] \) 上单调递增}
\end{problem}

\begin{answer}
B
\end{answer}

\begin{solution}
图象向右平移 \(\frac\pi2\) 后，函数变为
\[
y=3\sin\left(2\left(x-\frac\pi2\right)+\frac\pi3\right)=3\sin\left(2x-\frac{2\pi}{3}\right)
\]
其导数为
\[
y'=6\cos\left(2x-\frac{2\pi}{3}\right)
\]
当 \(x\in\left[\frac\pi{12},\frac{7\pi}{12}\right]\) 时，
\[
2x-\frac{2\pi}{3}\in\left[-\frac\pi2,\frac\pi2\right]
\]
从而 \(y'\geqslant0\)，函数在该区间上单调递增. 故选 B.
\end{solution}

\begin{problem}
已知点 \(A(-2,3)\) 在抛物线 \(C: y^{2} = 2px\) 的准线上，过点 \(A\) 的直线与 \(C\) 在第一象限相切于点 \(B\)，记 \(C\) 的焦点为 \(F\)，则直线 \(BF\) 的斜率为（\quad）

\choices
    {\(\frac{1}{2}\)}
    {\(\frac{2}{3}\)}
    {\(\frac{3}{4}\)}
    {\(\frac{4}{3}\)}
\end{problem}

\begin{answer}
D
\end{answer}

\begin{solution}
抛物线 \(y^2=2px\) 的准线为 \(x=-\frac p2\). 由点 \(A(-2,3)\) 在准线上，得 \(p=4\)，所以抛物线为 \(y^2=8x\)，焦点为 \(F(2,0)\).

设切点 \(B=(x_1,y_1)\)，其中 \(x_1>0\)，\(y_1>0\). 抛物线的切线方程为
\[
yy_1=4(x+x_1)
\]
由于点 \(A\) 在切线上，且 \(y_1^2=8x_1\)，有
\[
3y_1=4(x_1-2)=\frac{y_1^2}{2}-8
\]
所以
\[
y_1^2-6y_1-16=0
\]
即 \((y_1-8)(y_1+2)=0\). 由 \(y_1>0\)，得 \(y_1=8\)，进而 \(x_1=8\). 因此
\[
k_{BF}=\frac{8-0}{8-2}=\frac43
\]
故选 D.
\end{solution}

\begin{problem}
当 \(x \in [-2, 1]\) 时，不等式 \(ax^3 - x^2 + 4x + 3 \geqslant 0\) 恒成立，则实数 \(a\) 的取值范围是（\quad）

\choices
    {\([-5, - 3]\)}
    {\(\left[-6, -\frac{9}{8}\right]\)}
    {\([-6, - 2]\)}
    {\([-4, - 3]\)}
\end{problem}

\begin{answer}
C
\end{answer}

\begin{solution}
令
\[
F_a(x)=ax^3-x^2+4x+3
\]
由题意，取 \(x=-1\) 和 \(x=1\)，分别得到
\(-a-2=F_a(-1)\geqslant0\)，\(a+6=F_a(1)\geqslant0\).
所以必须有 \(a\leqslant-2\) 且 \(a\geqslant-6\)，即 \(a\in[-6,-2]\).

反之，若 \(a\in[-6,-2]\)，当 \(x\in[-2,0]\) 时，令 \(t=-x\)，则 \(t\in[0,2]\)，且
\[
F_a(x)=-at^3-t^2-4t+3
\geqslant2t^3-t^2-4t+3
=(t-1)^2(2t+3)\geqslant0
\]
当 \(x\in[0,1]\) 时，由 \(a\geqslant-6\) 得
\[
F_a(x)\geqslant-6x^3-x^2+4x+3
=(1-x)(6x^2+7x+3)\geqslant0
\]
因此不等式在整个区间 \([-2,1]\) 上恒成立的充要条件是 \(a\in[-6,-2]\)，故选 C.
\end{solution}

\begin{problem}
已知定义在 \([0,1]\) 上的函数 \(f(x)\) 满足：

\begin{circlelist}
\item \(f(0) = f(1) = 0\)；
\item 对所有 \(x\)，\(y\in[0,1]\)，且 \(x \neq y\)，有 \(|f(x) - f(y)| < \frac{1}{2} |x - y|\).
\end{circlelist}

若对所有 \(x\)，\(y\in[0,1]\)，\(|f(x) - f(y)| < k\) 恒成立，则 \(k\) 的最小值为（\quad）

\choices
    {\(\frac{1}{2}\)}
    {\( \frac{1}{4} \)}
    {\( \frac{1}{2\pi} \)}
    {\(\frac{1}{8}\)}
\end{problem}

\begin{answer}
B
\end{answer}

\begin{solution}
对任意 \(0\leqslant x<y\leqslant1\)，由题设有 \(|f(x)-f(y)|<\frac12(y-x)\). 当 \((x,y)\neq(0,1)\) 时，由 \(f(0)=f(1)=0\) 及题设，
\[
|f(x)-f(y)|\leqslant|f(x)-f(0)|+|f(1)-f(y)|<\frac12x+\frac12(1-y)=\frac12(1-y+x)
\]
因此
\[
|f(x)-f(y)|<\frac12\min\{y-x,1-y+x\}\leqslant\frac14
\]
若 \((x,y)=(0,1)\)，则 \(|f(x)-f(y)|=0<\frac14\)，所以结论对任意 \(x<y\) 都成立.
交换 \(x\)，\(y\) 后结论对任意两个点都成立，所以 \(k=\frac14\) 可以满足条件.

下面说明这个常数不能再减小. 对任意 \(0<\varepsilon<\frac12\)，令
\[
f_\varepsilon(x)=\left(\frac12-\varepsilon\right)\min\{x,1-x\}
\]
该函数满足 \(f_\varepsilon(0)=f_\varepsilon(1)=0\)，且其 Lipschitz 常数为 \(\frac12-\varepsilon<\frac12\)，所以满足题设条件. 但
\[
\left|f_\varepsilon\left(\frac12\right)-f_\varepsilon(0)\right|=\frac14-\frac\varepsilon2
\]
令 \(\varepsilon\) 充分接近 \(0\)，这个差值可任意接近 \(\frac14\)，故任何小于 \(\frac14\) 的常数都不能对所有这样的函数成立. 因此 \(k\) 的最小值为 \(\frac14\)，故选 B.
\end{solution}

\section{填空题}

\begin{problem}
执行如图的程序框图，若输入 \(x = 9\)，则输出 \(y =\)\fillinblank{}.
\bitmapfigure[max width=0.42\linewidth]{img/2014/liaoning_science/q13_fig1.png}
\end{problem}

\begin{answer}
\(\frac{29}{9}\)
\end{answer}

\begin{solution}
输入 \(x=9\) 后，第一次计算得到 \(y=\frac93+2=5\)，且 \(|y-x|=4\geqslant1\)，于是令 \(x=5\).

第二次计算得到 \(y=\frac53+2=\frac{11}{3}\)，且 \(|y-x|=\frac43\geqslant1\)，于是令 \(x=\frac{11}{3}\).

第三次计算得到
\(y=\frac{11/3}{3}+2=\frac{29}{9}\)，\(|y-x|=\left|\frac{29}{9}-\frac{11}{3}\right|=\frac49<1\).
此时满足判断条件，输出 \(y=\frac{29}{9}\).
\end{solution}

\begin{problem}
正方形的四个顶点 \(A(-1, -1)\)，\(B(1, -1)\)，\(C(1, 1)\)，\(D(-1, 1)\) 分别在抛物线 \(y = -x^2\) 和 \(y = x^2\) 上，如图所示，若将一个质点随机投入正方形 \(ABCD\) 中，则质点落在阴影区域的概率是\fillinblank{}.
\bitmapfigure[max width=0.44\linewidth]{img/2014/liaoning_science/q14_fig1.png}
\end{problem}

\begin{answer}
\(\frac{2}{3}\)
\end{answer}

\begin{solution}
正方形 \(ABCD\) 的面积为 \(4\). 阴影区域由上、下两部分组成，上部面积为
\[
\int_{-1}^{1}(1-x^2)\,\mathrm{d}x=\frac43
\]
下部与上部关于原点中心对称，面积也为 \(\frac43\). 因此阴影区域面积为 \(\frac83\)，所求概率为
\[
\frac{\frac83}{4}=\frac23
\]
\end{solution}

\begin{problem}
已知椭圆 \(C: \frac{x^2}{9} + \frac{y^2}{4} = 1\)，点 \(M\) 与 \(C\) 的焦点不重合，若 \(M\) 关于 \(C\) 的焦点的对称点分别为 \(A\)，\(B\)，线段 \(MN\) 的中点在 \(C\) 上，则 \(|AN| + |BN| =\fillinblank{}\).
\end{problem}

\begin{answer}
\(12\)
\end{answer}

\begin{solution}
椭圆 \(C\) 的长半轴为 \(3\). 设其两个焦点为 \(F_1\)，\(F_2\)，设线段 \(MN\) 的中点为 \(N_0\)，则 \(N_0\in C\). 由中心对称关系，
\(A=2F_1-M\)，\(B=2F_2-M\)，\(N=2N_0-M\).
从而
\(|AN|=|2F_1-2N_0|=2|F_1N_0|\)，\(|BN|=2|F_2N_0|\).
又椭圆的定义给出
\[
|F_1N_0|+|F_2N_0|=2\times3=6
\]
所以
\[
|AN|+|BN|=2\bigl(|F_1N_0|+|F_2N_0|\bigr)=12
\]
\end{solution}

\begin{problem}
对于 \(c > 0\)，当非零实数 \(a\)，\(b\) 满足 \(4a^2 - 2ab + 4b^2 - c = 0\) 且使 \(|2a + b|\) 最大时，\(\frac{3}{a} - \frac{4}{b} + \frac{5}{c}\) 的最小值为\fillinblank{}.
\end{problem}

\begin{answer}
\(-2\)
\end{answer}

\begin{solution}
有恒等式
\[
\frac85\left(4a^2-2ab+4b^2\right)-(2a+b)^2=\frac35(2a-3b)^2\geqslant0
\]
由题设得
\[
|2a+b|\leqslant\sqrt{\frac{8c}{5}}
\]
且等号成立当且仅当 \(2a=3b\). 因此在 \(|2a+b|\) 取得最大值时，可设
\(a=3t\)，\(b=2t\)，\(40t^2=c\)，\(t\neq0\).
此时
\[
\frac3a-\frac4b+\frac5c=\frac1t-\frac2t+\frac5{40t^2}=-\frac1t+\frac1{8t^2}
\]
对固定的 \(c\)，取正的 \(t\) 可使上式小于取负的 \(t\)，所以令 \(u=\frac1t>0\)，问题化为求
\[
\frac{u^2}{8}-u=\frac18(u-4)^2-2
\]
的最小值. 当 \(u=4\) 时取得；由于 \(c=40t^2\)，且 \(c\) 可取任意正数，\(u\) 确实可取任意正数. 故最小值为 \(-2\).
\end{solution}

\section{解答题}

\begin{problem}
在 \(\triangle ABC\) 中，内角 \(A\)，\(B\)，\(C\) 的对边分别为 \(a\)，\(b\)，\(c\)，且 \(a > c\)，已知 \(\overrightarrow{BA} \cdot \overrightarrow{BC} = 2\)，\(\cos B = \frac{1}{3}\)，\(b = 3\)，求：

\begin{enumerate}
\item \(a\) 和 \(c\) 的值；
\item \(\cos (B - C)\) 的值.
\end{enumerate}
\end{problem}

\begin{solution}
因为 \(\lvert\overrightarrow{BA}\rvert=c\)，\(\lvert\overrightarrow{BC}\rvert=a\)，所以
\[
\overrightarrow{BA}\cdot\overrightarrow{BC}=ac\cos B=\frac{ac}{3}=2
\]
从而 \(ac=6\). 又由余弦定理，
\[
b^2=a^2+c^2-2ac\cos B=a^2+c^2-4
\]
代入 \(b=3\)，得 \(a^2+c^2=13\). 因此
\[
(a+c)^2=a^2+c^2+2ac=25
\]
由于 \(a\)，\(c\) 是边长，\(a+c=5\)，所以 \(a\)，\(c\) 是方程
\[
t^2-5t+6=0
\]
的两个根，即 \(a\)，\(c\) 为 \(2\)，\(3\). 由 \(a>c\)，得
\(a=3\)，\(c=2\).

再由余弦定理，
\[
c^2=a^2+b^2-2ab\cos C
\]
所以
\(4=9+9-18\cos C\)，\(\cos C=\frac79\).
因为 \(B\)，\(C\) 是三角形内角，
\(\sin B=\frac{2\sqrt2}{3}\)，\(\sin C=\sqrt{1-\frac{49}{81}}=\frac{4\sqrt2}{9}\).
于是
\[
\cos(B-C)=\cos B\cos C+\sin B\sin C
=\frac13\cdot\frac79+\frac{2\sqrt2}{3}\cdot\frac{4\sqrt2}{9}
=\frac{23}{27}
\]
\end{solution}

\begin{problem}
一家面包房根据以往某种面包的销售记录，绘制了日销售量的频率分布直方图，如图所示. 将日销售量落入各组的频率视为概率，并假设每天的销售量相互独立.

\begin{enumerate}
\item 求在未来连续 3 天里，有连续 2 天的日销售量都不低于 100 个且另 1 天的日销售量低于 50 个的概率；
\item 用 \(X\) 表示在未来 3 天里日销售量不低于 100 个的天数，求随机变量 \(X\) 的分布列，期望 \(E(X)\) 及方差 \(D(X)\).
\end{enumerate}

\bitmapfigure[max width=0.42\linewidth]{img/2014/liaoning_science/q18_fig1.png}
\end{problem}

\begin{solution}
由直方图可知各组组距为 \(50\)，所以日销售量低于 \(50\) 个的概率为
\[
p_0=50\times0.003=0.15=\frac3{20}
\]
日销售量不低于 \(100\) 个的概率为
\[
p=50\times(0.006+0.004+0.002)=0.6=\frac35
\]

题目所述事件中，低于 \(50\) 个的那一天只能在连续三天的第一天或第三天，因此该事件的概率为
\[
2\times\frac3{20}\times\left(\frac35\right)^2
=\frac{27}{250}
\]

三天的销售量相互独立，所以 \(X\) 服从参数为 \(3\)，\(\frac35\) 的二项分布. 对 \(k=0\)，\(1\)，\(2\)，\(3\)，
\[
P(X=k)=\mathrm C_3^k\left(\frac35\right)^k\left(\frac25\right)^{3-k}
\]
因此 \(X\) 的分布列为
\begin{center}
\begin{tabular}{c|cccc}
\(X\) & \(0\) & \(1\) & \(2\) & \(3\) \\
\hline
\(P\) & \(\dfrac8{125}\) & \(\dfrac{36}{125}\) & \(\dfrac{54}{125}\) & \(\dfrac{27}{125}\)
\end{tabular}
\end{center}
并且
\(E(X)=3\times\frac35=\frac95\)，\(D(X)=3\times\frac35\times\frac25=\frac{18}{25}\).
\end{solution}

\begin{problem}
如图，\(\triangle ABC\) 和 \(\triangle BCD\) 所在平面互相垂直，且 \(AB = BC = BD = 2\)，\(\angle ABC = \angle DBC = 120^{\circ}\)，\(E\)，\(F\) 分别为 \(AC\)，\(DC\) 的中点.

\begin{enumerate}
\item 求证：\(EF \perp BC\)；
\item 求二面角 \(E - BF - C\) 的正弦值.
\end{enumerate}

\bitmapfigure[max width=0.34\linewidth]{img/2014/liaoning_science/q19_fig1.png}
\end{problem}

\begin{solution}
以 \(B\) 为原点，取 \(BC\) 所在直线为 \(x\) 轴，取三角形 \(ABC\) 所在平面为 \(xy\) 平面，三角形 \(BCD\) 所在平面为 \(xz\) 平面. 由题设可取
\(B=(0,0,0)\)，\(C=(2,0,0)\)，\(A=(-1,\sqrt3,0)\)，\(D=(-1,0,\sqrt3)\).
于是
\(E=\left(\frac12,\frac{\sqrt3}{2},0\right)\)，\(F=\left(\frac12,0,\frac{\sqrt3}{2}\right)\).
所以
\(\overrightarrow{EF}=\left(0,-\frac{\sqrt3}{2},\frac{\sqrt3}{2}\right)\)，\(\overrightarrow{BC}=(2,0,0)\)
从而 \(\overrightarrow{EF}\cdot\overrightarrow{BC}=0\)，故 \(EF\perp BC\).

设二面角 \(E-BF-C\) 为 \(\theta\). 平面 \(EBF\) 和 \(CBF\) 的法向量可分别取
\[
\bs{n}_1=\overrightarrow{BE}\times\overrightarrow{BF}
=\left(\frac34,-\frac{\sqrt3}{4},-\frac{\sqrt3}{4}\right)
\]
\[
\bs{n}_2=\overrightarrow{BC}\times\overrightarrow{BF}=(0,-\sqrt3,0)
\]
于是
\(\lvert\bs{n}_1\rvert=\frac{\sqrt{15}}4\)，\(\lvert\bs{n}_2\rvert=\sqrt3\)，\(\bs{n}_1\cdot\bs{n}_2=\frac34\).
因此
\[
\cos\theta=\frac{\bs{n}_1\cdot\bs{n}_2}{\lvert\bs{n}_1\rvert\lvert\bs{n}_2\rvert}=\frac1{\sqrt5}
\]
从而
\[
\sin\theta=\sqrt{1-\frac15}=\frac2{\sqrt5}
\]
\end{solution}

\begin{problem}
圆 \(x^{2} + y^{2} = 4\) 的切线与 \(x\) 轴正半轴，\(y\) 轴正半轴围成一个三角形，当该三角形面积最小时，切点为 \(P\)（如图）. 双曲线 \(C_{1}: \frac{x^{2}}{a^{2}} - \frac{y^{2}}{b^{2}} = 1\) 过点 \(P\) 且离心率为 \(\sqrt{3}\).
\begin{enumerate}
\item 求 \(C_1\) 的方程；
\item 椭圆 \(C_2\) 过点 \(P\) 且与 \(C_1\) 有相同的焦点，直线 \(l\) 过 \(C_2\) 的右焦点且与 \(C_2\) 交于 \(A\)，\(B\) 两点，若以线段 \(AB\) 为直径的圆过点 \(P\)，求 \(l\) 的方程.
\end{enumerate}

\FigureLayoutDeclare{img/2014/liaoning_science/q20_fig1.png}{5.6cm}{0bp 107bp 171bp 36bp}
\bitmapfigure[max width=0.34\linewidth]{img/2014/liaoning_science/q20_fig1.png}
\end{problem}

\begin{solution}
\begin{enumerate}
\item
设切点为 \(P=(x_0,y_0)\)，其中 \(x_0>0\)，\(y_0>0\). 圆的切线方程为
\[
x_0x+y_0y=4
\]
它在两坐标轴正半轴上的截距分别为 \(\frac4{x_0}\)，\(\frac4{y_0}\)，所以三角形面积为
\[
S=\frac12\cdot\frac4{x_0}\cdot\frac4{y_0}=\frac8{x_0y_0}
\]
又 \(x_0^2+y_0^2=4\)，故
\[
x_0y_0\leqslant\frac{x_0^2+y_0^2}{2}=2
\]
当且仅当 \(x_0=y_0=\sqrt2\) 时等号成立. 因而面积最小时
\[
P=(\sqrt2,\sqrt2)
\]

\item
双曲线 \(C_1\) 的离心率满足
\[
3=e^2=1+\frac{b^2}{a^2}
\]
所以 \(b^2=2a^2\). 由于 \(P\in C_1\)，
\[
\frac2{a^2}-\frac2{b^2}=1
\]
代入 \(b^2=2a^2\)，得 \(a^2=1\)，\(b^2=2\)，故
\(C_1:\)，\(x^2-\frac{y^2}{2}=1\).

双曲线 \(C_1\) 的焦点为 \((\pm\sqrt3,0)\). 设椭圆 \(C_2\) 的方程为
\[
\frac{x^2}{\alpha^2}+\frac{y^2}{\beta^2}=1
\]
它与 \(C_1\) 有相同焦点，所以 \(\alpha^2-\beta^2=3\). 又 \(P\in C_2\)，故
\[
\frac2{\alpha^2}+\frac2{\beta^2}=1
\]
联立两式，得 \(\alpha^2=6\)，\(\beta^2=3\)，即
\(C_2:\)，\(\frac{x^2}{6}+\frac{y^2}{3}=1\).

记 \(C_2\) 的右焦点为 \(F=(\sqrt3,0)\). 先设 \(l\) 不是竖直直线，写成
\[
y=k(x-\sqrt3)
\]
用参数 \(t\) 表示直线上的点
\[
X(t)=(\sqrt3+t,kt)
\]
将其代入 \(C_2\)，交点参数 \(t_1\)，\(t_2\) 是方程
\[
(1+2k^2)t^2+2\sqrt3t-3=0
\]
的两个根，因此
\(t_1+t_2=-\frac{2\sqrt3}{1+2k^2}\)，\(t_1t_2=-\frac3{1+2k^2}\).
以 \(AB\) 为直径的圆过 \(P\)，等价于
\[
(P-X(t_1))\cdot(P-X(t_2))=0
\]
令 \(Q=P-F=(\sqrt2-\sqrt3,\sqrt2)\)，则
\(\lvert Q\rvert^2=7-2\sqrt6\)，\(Q\cdot(1,k)=\sqrt2(1+k)-\sqrt3\).
利用上面的根和与根积，整理得
\[
(7-2\sqrt6)(1+2k^2)+2\sqrt3\bigl(\sqrt2(1+k)-\sqrt3\bigr)-3(1+k^2)=0
\]
即
\[
(11-4\sqrt6)k^2+2\sqrt6k-2=0
\]
解得
\[
k=\frac{2+3\sqrt6}{25}\quad\text{或}\quad k=-(2+\sqrt6)
\]
竖直直线 \(x=\sqrt3\) 与 \(C_2\) 的交点为 \(\left(\sqrt3,\sqrt{\frac32}\right)\)，\(\left(\sqrt3,-\sqrt{\frac32}\right)\)，此时
\[
(P-A)\cdot(P-B)=\left(\sqrt2-\sqrt3\right)^2+\left(\sqrt2-\sqrt{\frac32}\right)\left(\sqrt2+\sqrt{\frac32}\right)
=\frac{11}{2}-2\sqrt6\neq0
\]
所以竖直直线不是所求直线. 因此
\(l:\)，\(y=\frac{2+3\sqrt6}{25}(x-\sqrt3) \quad\text{或}\quad y=-(2+\sqrt6)(x-\sqrt3)\).
\end{enumerate}
\end{solution}

\begin{problem}
已知函数 \(f(x) = (\cos x - x)(\pi + 2x) - \frac{8}{3} (\sin x + 1)\)，\(g(x) = 3(x - \pi)\cos x - 4(1 + \sin x)\ln \left(3 - \frac{2x}{\pi}\right)\).

\begin{enumerate}
\item 证明：存在唯一 \(x_0 \in \left(0, \frac{\pi}{2}\right)\)，使 \(f(x_0) = 0\)；
\item 证明：存在唯一 \(x_{1} \in \left(\frac{\pi}{2}, \pi\right)\)，使 \(g(x_{1}) = 0\)，且对（1）中的 \(x_{0}\)，有 \(x_{0} + x_{1} < \pi\).
\end{enumerate}
\end{problem}

\begin{solution}
\begin{enumerate}
\item
在 \(\left(0,\frac\pi2\right)\) 上，
\[
f'(x)=-(\pi+2x)(1+\sin x)-2x-\frac23\cos x<0
\]
又
\(f(0)=\pi-\frac83>0\)，\(f\left(\frac\pi2\right)=-\pi^2-\frac{16}{3}<0\).
所以由零点存在定理，存在 \(x_0\in\left(0,\frac\pi2\right)\) 使 \(f(x_0)=0\)，并由 \(f\) 的严格单调性知该零点唯一.

\item
令 \(t=\pi-x\)，则 \(t\in\left(0,\frac\pi2\right)\)，并定义
\[
H(t)=\frac{3t\cos t}{4(1+\sin t)}-\ln\left(1+\frac{2t}{\pi}\right)
\]
直接代入可得
\[
g(\pi-t)=4(1+\sin t)H(t)
\]
由
\[
H'(t)=\frac{3(\cos t-t)}{4(1+\sin t)}-\frac2{\pi+2t}
=\frac{3f(t)}{4(\pi+2t)(1+\sin t)}
\]
在 \(\left(0,\frac\pi2\right)\) 上，分母为正，而 \(f\) 在前面已证明严格递减且唯一零点为 \(x_0\)，所以 \(H\) 在 \(\left(0,x_0\right)\) 上严格递增，在 \(\left(x_0,\frac\pi2\right)\) 上严格递减. 又
\(H(0)=0\)，\(H\left(\frac\pi2\right)=-\ln2<0\).
因此 \(H(t)>0\)（\(0<t\leqslant x_0\)），并由零点存在定理及严格单调性，存在唯一
\(t_1\in\left(x_0,\frac\pi2\right)\) 使 \(H(t_1)=0\). 由于
\(g(\pi-t)=4(1+\sin t)H(t)\)，且 \(1+\sin t>0\)，故
\[
x_1=\pi-t_1
\]
是 \(\left(\frac\pi2,\pi\right)\) 内 \(g(x)=0\) 的唯一解. 最后 \(t_1>x_0\)，所以
\[
x_0+x_1=x_0+\pi-t_1<\pi
\]
\end{enumerate}
\end{solution}

\section{选考题：解答题}

\begin{problem}
如图，\(EP\) 交圆于 \(E\)，\(C\) 两点，\(PD\) 切圆于 \(D\)，\(G\) 为 \(CE\) 上一点且 \(PG = PD\)，连接 \(DG\) 并延长交圆于点 \(A\)，作弦 \(AB\) 垂直 \(EP\)，垂足为 \(F\).
\begin{enumerate}
\item 求证：\(AB\) 为圆的直径；
\item 若 \(AC = BD\)，求证：\(AB = ED\).
\end{enumerate}

\FigureLayoutDeclare{img/2014/liaoning_science/q22_fig1.png}{7.4cm}{123bp 66bp 171bp 45bp}
\bitmapfigure[max width=0.48\linewidth]{img/2014/liaoning_science/q22_fig1.png}
\end{problem}

\begin{solution}
\begin{enumerate}
\item
设圆心为 \(O\)，半径为 \(r\). 建立直角坐标系，使 \(O=(0,0)\)，\(EP\) 为直线 \(y=h\). 按图取
\(E=(-u,h)\)，\(C=(u,h)\)，\(P=(p,h)\)
其中 \(-r<h<0\)，\(u^2+h^2=r^2\)，\(p>u>0\). 设 \(G=(g,h)\). 由 \(G\) 在线段 \(CE\) 上以及 \(PG=PD\) 和点 \(P\) 的切线长公式，记
\(t=PG=PD=\sqrt{p^2-u^2}\)，\(g=p-t\).
令 \(A_0=(0,-r)\)，并令直线 \(A_0G\) 与圆除 \(A_0\) 外的另一个交点为 \(D_0\). 以参数 \(\lambda\) 表示直线 \(A_0G\) 上的点：
\[
X(\lambda)=\bigl(\lambda g,-r+\lambda(h+r)\bigr)
\]
将其代入圆方程，除去 \(\lambda=0\) 后得到
\(\lambda_0=\frac{2r(h+r)}{g^2+(h+r)^2}\)，\(D_0=X(\lambda_0)\).
由 \(g=p-t\) 及 \(t^2=p^2-u^2\)，有
\[
2\bigl(gp+(h+r)h\bigr)-\bigl(g^2+(h+r)^2\bigr)
=p^2-t^2+h^2-r^2=0
\]
所以
\[
D_0\cdot P=-rh+\lambda_0\bigl(gp+(h+r)h\bigr)
=-rh+r(h+r)=r^2
\]
由于 \(t<p\) 且 \(t>p-u\)，有 \(0<g<u\). 又

\[
2r(h+r)-\bigl(g^2+(h+r)^2\bigr)=u^2-g^2>0
\]

故 \(\lambda_0>1\)，从而 \(D_0\) 位于 \(y=h\) 的上方且横坐标为正. 由于

\(\lvert D_0\rvert=r\)，\(D_0\cdot P=r^2\)

有 \((P-D_0)\cdot D_0=0\)，所以 \(PD_0\) 是圆在 \(D_0\) 处的切线. 这正是图中位于弦 \(EC\) 上方的切点 \(D\)，故 \(D_0=D\)，从而 \(A_0\)，\(D\)，\(G\) 共线. 因而题设中的点 \(A=A_0=(0,-r)\).

直线 \(AB\perp EP\)，所以 \(AB\) 是过 \(A\) 的竖直弦，另一交点为 \(B=(0,r)\). 因此 \(O\) 是 \(AB\) 的中点，\(AB\) 为圆的直径.

\item
在此基础上，令
\(E=(-r\cos\theta,-r\sin\theta)\)，\(C=(r\cos\theta,-r\sin\theta)\)
\(D=(r\cos\varphi,r\sin\varphi)\)，\(0<\theta\)，\(\varphi<\frac\pi2\).
此时 \(A=(0,-r)\)，\(B=(0,r)\)，所以
\(AC^2=2r^2(1-\sin\theta)\)，\(BD^2=2r^2(1-\sin\varphi)\).
若 \(AC=BD\)，则 \(\sin\theta=\sin\varphi\)，从而 \(\theta=\varphi\). 于是 \(D=-E\)，即 \(ED\) 也是圆的直径，故
\[
ED=2r=AB
\]
\end{enumerate}
\end{solution}

\begin{problem}
将圆 \(x^{2} + y^{2} = 1\) 上每一点的横坐标保持不变，纵坐标变为原来的 2 倍，得曲线 \(C\).

\begin{enumerate}
\item 写出 \(C\) 的参数方程；
\item 设直线 \(l: 2x + y - 2 = 0\) 与 \(C\) 的交点为 \(P_{1}\)，\(P_{2}\)，以坐标原点为极点，\(x\) 轴正半轴为极轴建立极坐标系. 求过线段 \(P_{1}P_{2}\) 的中点且与 \(l\) 垂直的直线的极坐标方程.
\end{enumerate}
\end{problem}

\begin{solution}
\begin{enumerate}
\item
原圆的参数方程可写为
\(x=\cos\theta\)，\(y=\sin\theta\)，\(0\leqslant\theta<2\pi\).
变换后横坐标不变，纵坐标变为原来的 \(2\) 倍，所以曲线 \(C\) 的参数方程为
\[
\begin{cases}
x=\cos\theta,\\
y=2\sin\theta,
\end{cases}
\qquad 0\leqslant\theta<2\pi
\]

\item
消去参数可得 \(C\) 的方程为
\[
x^2+\frac{y^2}{4}=1
\]
直线 \(l\) 与 \(C\) 联立，得
\(y=2-2x\)，\(x^2+(1-x)^2=1\).
所以 \(x=0\) 或 \(x=1\)，交点为
\(P_1=(0,2)\)，\(P_2=(1,0)\).
线段 \(P_1P_2\) 的中点为
\[
M=\left(\frac12,1\right)
\]
直线 \(l\) 的斜率为 \(-2\)，故过 \(M\) 且与 \(l\) 垂直的直线斜率为 \(\frac12\)，其直角坐标方程为
\(y-1=\frac12\left(x-\frac12\right)\)，\(y=\frac12x+\frac34\).
在极坐标系中代入 \(x=\rho\cos\theta\)，\(y=\rho\sin\theta\)，得到所求极坐标方程
\[
\rho\left(2\sin\theta-\cos\theta\right)=\frac32
\]
\end{enumerate}
\end{solution}

\begin{problem}
设函数 \( f(x) = 2|x - 1| + x - 1 \)，\( g(x) = 16x^{2} - 8x + 1 \)，记 \( f(x) \leqslant 1 \) 的解集为 \(M\)，\( g(x) \leqslant 4 \) 的解集为 \(N\).

\begin{enumerate}
\item 求 \(M\)；
\item 当 \(x \in M \cap N\) 时，证明：\(x^{2}f(x) + x[f(x)]^{2} \leqslant \frac{1}{4}\).
\end{enumerate}
\end{problem}

\begin{solution}
\begin{enumerate}
\item
当 \(x<1\) 时，\(f(x)=1-x\)，所以 \(f(x)\leqslant1\) 等价于 \(x\geqslant0\). 当 \(x\geqslant1\) 时，\(f(x)=3x-3\)，所以 \(f(x)\leqslant1\) 等价于 \(x\leqslant\frac43\). 因此
\[
M=[0,\tfrac43]
\]

又
\[
g(x)=(4x-1)^2\leqslant4
\Longleftrightarrow -2\leqslant4x-1\leqslant2
\]
所以
\(N=\left[-\frac14,\frac34\right]\)，\(M\cap N=\left[0,\frac34\right]\).
\item
当 \(x\in M\cap N\) 时，\(x\leqslant1\)，故 \(f(x)=1-x\). 从而
\[
x^2f(x)+x[f(x)]^2
=x^2(1-x)+x(1-x)^2
=x(1-x)
\leqslant\frac14
\]
其中最后一个不等号由 \(\left(x-\frac12\right)^2\geqslant0\) 得到. 结论成立.
\end{enumerate}
\end{solution}
